\documentclass[11pt]{article}
\usepackage[margin=1in]{geometry}
\usepackage{amsmath}
\usepackage{amssymb}
\usepackage{enumitem}
\usepackage{array}
\usepackage{xcolor}

\setlength{\parindent}{0pt}
\setlength{\parskip}{4pt}

\begin{document}

\begin{center}
  {\LARGE\bfseries Fractions: Worked Solutions}\\[4pt]
  {\large Methods and answers}
\end{center}

%================================================================
\section*{Simplifying fractions}
In each fraction, divide the numerator and denominator by their highest common factor.

\begin{enumerate}[label=\arabic*., start=1, leftmargin=*, itemsep=6pt]
\item Simplify \(\frac{4}{8}\).
\textbf{Worked solution.}
\[
\frac{4}{8}=\frac{4\div4}{8\div4}=\frac12.
\]

\item Simplify \(\frac{6}{9}\).
\textbf{Worked solution.}
\[
\frac{6}{9}=\frac{6\div3}{9\div3}=\frac23.
\]

\item Simplify \(\frac{12}{20}\).
\textbf{Worked solution.}
\[
\frac{12}{20}=\frac{12\div4}{20\div4}=\frac35.
\]

\item Simplify \(\frac{24}{36}\).
\textbf{Worked solution.}
\[
\frac{24}{36}=\frac{24\div12}{36\div12}=\frac23.
\]

\item Simplify \(\frac{45}{60}\).
\textbf{Worked solution.}
\[
\frac{45}{60}=\frac{45\div15}{60\div15}=\frac34.
\]

\item Simplify \(\frac{42}{70}\).
\textbf{Worked solution.}
\[
\frac{42}{70}=\frac{42\div14}{70\div14}=\frac35.
\]

\item Simplify \(\frac{84}{144}\).
\textbf{Worked solution.}
\[
\frac{84}{144}=\frac{84\div12}{144\div12}=\frac{7}{12}.
\]

\item Simplify \(\frac{91}{117}\).
\textbf{Worked solution.}
\[
91=7\times13,\qquad 117=9\times13,
\]
so
\[
\frac{91}{117}=\frac{91\div13}{117\div13}=\frac79.
\]

\item Simplify \(\frac{315}{420}\).
\textbf{Worked solution.}
\[
\frac{315}{420}=\frac{315\div105}{420\div105}=\frac34.
\]

\item Simplify \(\frac{1001}{1309}\).
\textbf{Worked solution.}
Using the hint, both numbers are multiples of \(7\) and \(11\), so divide by \(77\):
\[
1001\div77=13,\qquad 1309\div77=17.
\]
Hence
\[
\frac{1001}{1309}=\frac{13}{17}.
\]
\end{enumerate}

%================================================================
\section*{Ordering fractions}
Rewrite each list using a common denominator, then compare the numerators.

\begin{enumerate}[label=\arabic*., start=1, leftmargin=*, itemsep=6pt]
\item Put \(\frac14,\ \frac12,\ \frac18\) in order, smallest first.
\textbf{Worked solution.}
\[
\frac14=\frac28,\qquad \frac12=\frac48,\qquad \frac18=\frac18.
\]
Since \(\frac18<\frac28<\frac48\),
\[
\boxed{\frac18<\frac14<\frac12}.
\]

\item Put \(\frac37,\ \frac67,\ \frac27\) in order, smallest first.
\textbf{Worked solution.}
They already have denominator \(7\):
\[
\boxed{\frac27<\frac37<\frac67}.
\]

\item Put \(\frac12,\ \frac34,\ \frac58\) in order, smallest first.
\textbf{Worked solution.}
\[
\frac12=\frac48,\qquad \frac34=\frac68,\qquad \frac58=\frac58.
\]
Hence
\[
\boxed{\frac12<\frac58<\frac34}.
\]

\item Put \(\frac23,\ \frac56,\ \frac34\) in order, smallest first.
\textbf{Worked solution.}
\[
\frac23=\frac{8}{12},\qquad \frac56=\frac{10}{12},\qquad \frac34=\frac{9}{12}.
\]
Hence
\[
\boxed{\frac23<\frac34<\frac56}.
\]

\item Put \(\frac35,\ \frac{7}{10},\ \frac{13}{20}\) in order, smallest first.
\textbf{Worked solution.}
\[
\frac35=\frac{12}{20},\qquad \frac{7}{10}=\frac{14}{20},\qquad \frac{13}{20}=\frac{13}{20}.
\]
Hence
\[
\boxed{\frac35<\frac{13}{20}<\frac{7}{10}}.
\]

\item Put \(\frac59,\ \frac23,\ \frac{7}{12}\) in order, smallest first.
\textbf{Worked solution.}
\[
\frac59=\frac{20}{36},\qquad \frac23=\frac{24}{36},\qquad \frac{7}{12}=\frac{21}{36}.
\]
Hence
\[
\boxed{\frac59<\frac{7}{12}<\frac23}.
\]

\item Put \(\frac45,\ \frac79,\ \frac{11}{15}\) in order, \textbf{largest} first.
\textbf{Worked solution.}
\[
\frac45=\frac{36}{45},\qquad \frac79=\frac{35}{45},\qquad \frac{11}{15}=\frac{33}{45}.
\]
Largest first:
\[
\boxed{\frac45>\frac79>\frac{11}{15}}.
\]

\item Put \(\frac58,\ \frac{7}{11},\ \frac{9}{14}\) in order, smallest first.
\textbf{Worked solution.}
Use denominator \(616\):
\[
\frac58=\frac{385}{616},\qquad \frac{7}{11}=\frac{392}{616},\qquad \frac{9}{14}=\frac{396}{616}.
\]
Hence
\[
\boxed{\frac58<\frac{7}{11}<\frac{9}{14}}.
\]

\item Which is closer to \(1\): \(\frac89\) or \(\frac{11}{12}\)?
\textbf{Worked solution.}
\[
1-\frac89=\frac19,\qquad 1-\frac{11}{12}=\frac{1}{12}.
\]
Since \(\frac{1}{12}<\frac19\), \(\frac{11}{12}\) is closer to \(1\).
\[
\boxed{\frac{11}{12}}
\]

\item Write a fraction that lies between \(\frac58\) and \(\frac23\).
\textbf{Worked solution.}
\[
\frac58=\frac{30}{48},\qquad \frac23=\frac{32}{48}.
\]
So one possible fraction is
\[
\boxed{\frac{31}{48}}.
\]
There are infinitely many correct answers; for example \(\frac{31}{48}\) works.
\end{enumerate}

%================================================================
\section*{Mixed numbers and improper fractions}
Convert by multiplying the whole number by the denominator and adding the numerator.

\begin{enumerate}[label=\arabic*., start=1, leftmargin=*, itemsep=6pt]
\item Write \(1\frac12\) as an improper fraction.
\textbf{Worked solution.}
\[
1\frac12=\frac{1\times2+1}{2}=\boxed{\frac32}.
\]

\item Write \(2\frac34\) as an improper fraction.
\textbf{Worked solution.}
\[
2\frac34=\frac{2\times4+3}{4}=\boxed{\frac{11}{4}}.
\]

\item Write \(\frac72\) as a mixed number.
\textbf{Worked solution.}
\[
7\div2=3\text{ remainder }1,
\]
so
\[
\frac72=\boxed{3\frac12}.
\]

\item Write \(4\frac25\) as an improper fraction.
\textbf{Worked solution.}
\[
4\frac25=\frac{4\times5+2}{5}=\boxed{\frac{22}{5}}.
\]

\item Write \(\frac{23}{4}\) as a mixed number.
\textbf{Worked solution.}
\[
23\div4=5\text{ remainder }3,
\]
so
\[
\frac{23}{4}=\boxed{5\frac34}.
\]

\item Write \(5\frac78\) as an improper fraction.
\textbf{Worked solution.}
\[
5\frac78=\frac{5\times8+7}{8}=\boxed{\frac{47}{8}}.
\]

\item Write \(\frac{100}{7}\) as a mixed number.
\textbf{Worked solution.}
\[
100\div7=14\text{ remainder }2,
\]
so
\[
\frac{100}{7}=\boxed{14\frac27}.
\]

\item Write \(12\frac59\) as an improper fraction.
\textbf{Worked solution.}
\[
12\frac59=\frac{12\times9+5}{9}=\boxed{\frac{113}{9}}.
\]

\item Write \(-3\frac25\) as an improper fraction.
\textbf{Worked solution.}
\[
-3\frac25=-\frac{3\times5+2}{5}=\boxed{-\frac{17}{5}}.
\]

\item The mixed number \(6\frac{a}{7}\) is the same as \(\frac{45}{7}\). Find \(a\).
\textbf{Worked solution.}
\[
6\frac{a}{7}=\frac{6\times7+a}{7}=\frac{42+a}{7}.
\]
So
\[
\frac{42+a}{7}=\frac{45}{7}
\]
and therefore
\[
42+a=45 \implies \boxed{a=3}.
\]
\end{enumerate}

%================================================================
\section*{Adding and subtracting fractions}
Use a common denominator before adding or subtracting.

\begin{enumerate}[label=\arabic*., start=1, leftmargin=*, itemsep=6pt]
\item \(\frac12+\frac14\).
\textbf{Worked solution.}
\[
\frac12+\frac14=\frac24+\frac14=\boxed{\frac34}.
\]

\item \(\frac23+\frac16\).
\textbf{Worked solution.}
\[
\frac23+\frac16=\frac46+\frac16=\boxed{\frac56}.
\]

\item \(\frac34-\frac18\).
\textbf{Worked solution.}
\[
\frac34-\frac18=\frac68-\frac18=\boxed{\frac58}.
\]

\item \(\frac13+\frac14\).
\textbf{Worked solution.}
\[
\frac13+\frac14=\frac{4}{12}+\frac{3}{12}=\boxed{\frac{7}{12}}.
\]

\item \(\frac35-\frac14\).
\textbf{Worked solution.}
\[
\frac35-\frac14=\frac{12}{20}-\frac{5}{20}=\boxed{\frac{7}{20}}.
\]

\item \(\frac56+\frac38\).
\textbf{Worked solution.}
\[
\frac56+\frac38=\frac{20}{24}+\frac{9}{24}=\frac{29}{24}=\boxed{1\frac{5}{24}}.
\]

\item \(\frac{7}{10}-\frac{2}{15}\).
\textbf{Worked solution.}
\[
\frac{7}{10}-\frac{2}{15}=\frac{21}{30}-\frac{4}{30}=\boxed{\frac{17}{30}}.
\]

\item \(\frac23+\frac34+\frac16\).
\textbf{Worked solution.}
\[
\frac23+\frac34+\frac16=\frac{8}{12}+\frac{9}{12}+\frac{2}{12}
=\frac{19}{12}=\boxed{1\frac{7}{12}}.
\]

\item \(\frac{5}{12}-\frac{7}{18}\).
\textbf{Worked solution.}
\[
\frac{5}{12}-\frac{7}{18}=\frac{15}{36}-\frac{14}{36}=\boxed{\frac{1}{36}}.
\]

\item What must be added to \(\frac59\) to make \(\frac78\)?
\textbf{Worked solution.}
We need
\[
\frac78-\frac59=\frac{63}{72}-\frac{40}{72}=\boxed{\frac{23}{72}}.
\]
\end{enumerate}

%================================================================
\section*{Converting fractions, decimals and percentages}
Remember: percent means \emph{out of 100}.

\begin{enumerate}[label=\arabic*., start=1, leftmargin=*, itemsep=6pt]
\item Write \(\frac12\) as a decimal and as a percentage.
\textbf{Worked solution.}
\[
\frac12=0.5=50\%.
\]

\item Write \(0.25\) as a fraction and as a percentage.
\textbf{Worked solution.}
\[
0.25=\frac{25}{100}=\frac14,\qquad 0.25=25\%.
\]

\item Write \(30\%\) as a decimal and as a fraction in simplest form.
\textbf{Worked solution.}
\[
30\%=\frac{30}{100}=0.3.
\]
As a fraction:
\[
\frac{30}{100}=\frac{3}{10}.
\]

\item Write \(\frac34\) as a percentage, and \(0.6\) as a percentage.
\textbf{Worked solution.}
\[
\frac34=\frac{75}{100}=75\%,\qquad 0.6=60\%.
\]

\item Write \(\frac{7}{20}\) as a decimal and as a percentage.
\textbf{Worked solution.}
\[
\frac{7}{20}=\frac{35}{100}=0.35=35\%.
\]

\item Write \(0.08\) as a percentage and as a fraction in simplest form.
\textbf{Worked solution.}
\[
0.08=8\%,\qquad 0.08=\frac{8}{100}=\frac{2}{25}.
\]

\item Write \(\frac38\) as a decimal and as a percentage.
\textbf{Worked solution.}
\[
\frac38=0.375,\qquad \frac38=37.5\%.
\]

\item Write \(12.5\%\) as a fraction in simplest form.
\textbf{Worked solution.}
\[
12.5\%=\frac{12.5}{100}=\frac{125}{1000}=\frac18.
\]

\item Write \(\frac56\) as a decimal, and as a percentage to \(1\) decimal place.
\textbf{Worked solution.}
\[
\frac56=0.8333\ldots
\]
As a percentage:
\[
\frac56\times100=83.333\ldots\%,
\]
so to \(1\) decimal place,
\[
\frac56=83.3\%.
\]

\item Write \(137.5\%\) as a mixed number in simplest form.
\textbf{Worked solution.}
\[
137.5\%=\frac{137.5}{100}=\frac{1375}{1000}=\frac{11}{8}=1\frac38.
\]
\end{enumerate}

%================================================================
\section*{Ordering fractions, decimals and percentages}
Convert everything to the same representation before ordering.

\begin{enumerate}[label=\arabic*., start=1, leftmargin=*, itemsep=6pt]
\item Put \(0.5,\ \frac14,\ 60\%\) in order, smallest first.
\textbf{Worked solution.}
\[
0.5=50\%,\qquad \frac14=25\%,\qquad 60\%=60\%.
\]
Hence
\[
\boxed{\frac14<0.5<60\%}.
\]

\item Put \(0.7,\ \frac34,\ 70\%\) in order, smallest first.
\textbf{Worked solution.}
\[
0.7=70\%,\qquad \frac34=75\%.
\]
So
\[
\boxed{0.7=70\%<\frac34}.
\]

\item Put \(\frac25,\ 45\%,\ 0.35\) in order, smallest first.
\textbf{Worked solution.}
\[
\frac25=0.4=40\%,\qquad 45\%=0.45,\qquad 0.35=35\%.
\]
Hence
\[
\boxed{0.35<\frac25<45\%}.
\]

\item Put \(0.62,\ \frac58,\ 63\%\) in order, smallest first.
\textbf{Worked solution.}
\[
0.62=62\%,\qquad \frac58=62.5\%,\qquad 63\%=63\%.
\]
Hence
\[
\boxed{0.62<\frac58<63\%}.
\]

\item Put \(\frac78,\ 0.87,\ 88\%\) in order, \textbf{largest} first.
\textbf{Worked solution.}
\[
\frac78=0.875=87.5\%,\qquad 0.87=87\%,\qquad 88\%=0.88.
\]
Largest first:
\[
\boxed{88\%>\frac78>0.87}.
\]

\item Put \(\frac13,\ 0.3,\ 33\%\) in order, smallest first.
\textbf{Worked solution.}
\[
\frac13=33.333\ldots\%,\qquad 0.3=30\%,\qquad 33\%=33\%.
\]
Hence
\[
\boxed{0.3<33\%<\frac13}.
\]

\item Put \(0.045,\ 4.5\%,\ \frac{1}{20}\) in order, smallest first.
\textbf{Worked solution.}
\[
4.5\%=0.045,\qquad \frac{1}{20}=0.05=5\%.
\]
Hence
\[
\boxed{0.045=4.5\%<\frac{1}{20}}.
\]

\item Put \(\frac56,\ 83\%,\ 0.835\) in order, smallest first.
\textbf{Worked solution.}
\[
\frac56=0.8333\ldots,\qquad 83\%=0.83,\qquad 0.835=0.835.
\]
Hence
\[
\boxed{83\%<\frac56<0.835}.
\]

\item Put \(\frac27,\ 0.28,\ 29\%\) in order, smallest first.
\textbf{Worked solution.}
\[
\frac27=0.285714\ldots,\qquad 0.28=28\%,\qquad 29\%=0.29.
\]
Hence
\[
\boxed{0.28<\frac27<29\%}.
\]

\item Put \(\frac38,\ 0.38,\ 37.5\%,\ \frac37\) in order, smallest first.
\textbf{Worked solution.}
\[
\frac38=0.375,\qquad 37.5\%=0.375,\qquad 0.38=0.38,\qquad \frac37\approx0.4286.
\]
Hence
\[
\boxed{\frac38=37.5\%<0.38<\frac37}.
\]

\emph{What do you notice?} \(37.5\%\) is exactly equal to \(\frac38\), while \(0.38\) is slightly larger than \(\frac38\). Also, \(\frac37\) is the largest of these values.
\end{enumerate}

%================================================================
\section*{Fractions of amounts (no calculator)}
Find the unit fraction first, then multiply.

\begin{enumerate}[label=\arabic*., start=1, leftmargin=*, itemsep=6pt]
\item \(\frac12\) of \(40\).
\textbf{Worked solution.}
\[
40\div2=\boxed{20}.
\]

\item \(\frac14\) of \(60\).
\textbf{Worked solution.}
\[
60\div4=\boxed{15}.
\]

\item \(\frac15\) of \(45\).
\textbf{Worked solution.}
\[
45\div5=\boxed{9}.
\]

\item \(\frac34\) of \(40\).
\textbf{Worked solution.}
\[
40\div4=10,\qquad 3\times10=\boxed{30}.
\]

\item \(\frac23\) of \(36\).
\textbf{Worked solution.}
\[
36\div3=12,\qquad 2\times12=\boxed{24}.
\]

\item \(\frac35\) of \(45\).
\textbf{Worked solution.}
\[
45\div5=9,\qquad 3\times9=\boxed{27}.
\]

\item \(\frac58\) of \(72\).
\textbf{Worked solution.}
\[
72\div8=9,\qquad 5\times9=\boxed{45}.
\]

\item \(\frac27\) of \(350\).
\textbf{Worked solution.}
\[
350\div7=50,\qquad 2\times50=\boxed{100}.
\]

\item \(\frac56\) of a number is \(45\). What is the number?
\textbf{Worked solution.}
If \(\frac56\) of the number is \(45\), then \(\frac16\) of the number is
\[
45\div5=9.
\]
So the whole number is
\[
9\times6=\boxed{54}.
\]

\item \(\frac38\) of \(\frac25\) of \(200\).
\textbf{Worked solution.}
\[
\frac25\text{ of }200=200\div5\times2=80.
\]
Then
\[
\frac38\text{ of }80=80\div8\times3=30.
\]
So
\[
\boxed{30}.
\]
\end{enumerate}

%================================================================
\section*{Fractions of amounts (calculator)}
The unit-fraction method also works with a calculator.

\begin{enumerate}[label=\arabic*., start=1, leftmargin=*, itemsep=6pt]
\item \(\frac14\) of \(386\).
\textbf{Worked solution.}
\[
386\div4=96.5.
\]

\item \(\frac35\) of \(245\).
\textbf{Worked solution.}
\[
245\div5=49,\qquad 49\times3=147.
\]

\item \(\frac27\) of \(483\).
\textbf{Worked solution.}
\[
483\div7=69,\qquad 69\times2=138.
\]

\item \(\frac{3}{16}\) of \(1024\).
\textbf{Worked solution.}
\[
1024\div16=64,\qquad 64\times3=192.
\]

\item \(\frac59\) of \(738\).
\textbf{Worked solution.}
\[
738\div9=82,\qquad 82\times5=410.
\]

\item \(\frac78\) of \pounds 312.40.
\textbf{Worked solution.}
\[
312.40\div8=39.05,\qquad 39.05\times7=273.35.
\]
So the answer is
\[
\boxed{\pounds273.35}.
\]

\item \(\frac{11}{13}\) of \(2405\).
\textbf{Worked solution.}
\[
2405\div13=185,\qquad 185\times11=2035.
\]

\item \(\frac47\) of \(3.5\) kg. Give your answer in grams.
\textbf{Worked solution.}
\[
3.5\text{ kg}=3500\text{ g}.
\]
\[
3500\div7=500,\qquad 500\times4=2000.
\]
So the answer is
\[
\boxed{2000\text{ g}}.
\]

\item \(\frac58\) of a number is \(215\). What is the number?
\textbf{Worked solution.}
If \(\frac58\) of the number is \(215\), then \(\frac18\) of the number is
\[
215\div5=43.
\]
So the whole number is
\[
43\times8=\boxed{344}.
\]

\item \(\frac{7}{12}\) of \(4\) hours \(48\) minutes. Give your answer in hours and minutes.
\textbf{Worked solution.}
\[
4\text{ h }48\text{ min}=288\text{ min}.
\]
\[
288\div12=24,\qquad 24\times7=168\text{ min}.
\]
\[
168\text{ min}=2\text{ h }48\text{ min}.
\]
So the answer is
\[
\boxed{2\text{ hours }48\text{ minutes}}.
\]
\end{enumerate}

%================================================================
\section*{Four operations with mixed numbers and negatives}
Convert mixed numbers to improper fractions before calculating.

\begin{enumerate}[label=\arabic*., start=1, leftmargin=*, itemsep=6pt]
\item \(1\frac12+2\frac14\).
\textbf{Worked solution.}
\[
1\frac12+2\frac14=\frac32+\frac94=\frac64+\frac94=\frac{15}{4}
=\boxed{3\frac34}.
\]

\item \(3\frac45-1\frac25\).
\textbf{Worked solution.}
\[
3\frac45-1\frac25=\frac{19}{5}-\frac75=\frac{12}{5}
=\boxed{2\frac25}.
\]

\item \(2\frac12+\left(-1\frac14\right)\).
\textbf{Worked solution.}
Adding a negative is the same as subtracting:
\[
2\frac12-1\frac14=\frac52-\frac54=\frac{10}{4}-\frac54
=\frac54=\boxed{1\frac14}.
\]

\item \(4\frac16-2\frac34\).
\textbf{Worked solution.}
\[
4\frac16-2\frac34=\frac{25}{6}-\frac{11}{4}
=\frac{50}{12}-\frac{33}{12}
=\frac{17}{12}
=\boxed{1\frac{5}{12}}.
\]

\item \(2\frac12\times4\).
\textbf{Worked solution.}
\[
2\frac12\times4=\frac52\times\frac41=\frac{20}{2}=\boxed{10}.
\]

\item \(-3\frac13+1\frac56\).
\textbf{Worked solution.}
\[
-3\frac13+1\frac56=-\frac{10}{3}+\frac{11}{6}
=-\frac{20}{6}+\frac{11}{6}
=-\frac96
=-\frac32
=\boxed{-1\frac12}.
\]

\item \(1\frac12\times2\frac23\).
\textbf{Worked solution.}
\[
1\frac12\times2\frac23=\frac32\times\frac83
=\frac{24}{6}
=\boxed{4}.
\]

\item \(3\frac34\div1\frac12\).
\textbf{Worked solution.}
\[
3\frac34\div1\frac12=\frac{15}{4}\div\frac32
=\frac{15}{4}\times\frac23
=\frac{30}{12}
=\frac52
=\boxed{2\frac12}.
\]

\item \(-5\frac13\div\frac23\).
\textbf{Worked solution.}
\[
-5\frac13\div\frac23=-\frac{16}{3}\times\frac32
=-\frac{48}{6}
=\boxed{-8}.
\]

\item \(-2\frac14\times\left(-1\frac13\right)+1\frac12\).
\textbf{Worked solution.}
\[
-2\frac14=-\frac94,\qquad -1\frac13=-\frac43.
\]
First multiply:
\[
-\frac94\times-\frac43=\frac{36}{12}=3.
\]
Then add:
\[
3+1\frac12=3+\frac32=\frac92=\boxed{4\frac12}.
\]
\end{enumerate}

%================================================================
\section*{Word problems with mixed numbers}

\begin{enumerate}[label=\arabic*., start=1, leftmargin=*, itemsep=6pt]
\item A recipe needs \(1\frac12\) cups of flour. Alex makes a double batch. How much flour is needed?
\textbf{Worked solution.}
Double means multiply by \(2\):
\[
2\times1\frac12=2\times\frac32=3.
\]
So Alex needs
\[
\boxed{3\text{ cups}}.
\]

\item A plank is \(3\frac14\) m long. A piece \(1\frac12\) m long is cut off. How much plank is left?
\textbf{Worked solution.}
\[
3\frac14-1\frac12=\frac{13}{4}-\frac64=\frac74=1\frac34.
\]
So
\[
\boxed{1\frac34\text{ m}}.
\]

\item Priya has two ribbons, \(2\frac13\) m and \(1\frac56\) m long. What is their total length?
\textbf{Worked solution.}
\[
2\frac13+1\frac56=\frac73+\frac{11}{6}
=\frac{14}{6}+\frac{11}{6}
=\frac{25}{6}
=4\frac16.
\]
So
\[
\boxed{4\frac16\text{ m}}.
\]

\item Sam runs \(2\frac34\) miles on Monday and \(3\frac12\) miles on Tuesday. How much further did he run on Tuesday?
\textbf{Worked solution.}
\[
3\frac12-2\frac34=\frac72-\frac{11}{4}
=\frac{14}{4}-\frac{11}{4}
=\frac34.
\]
So Sam ran
\[
\boxed{\frac34\text{ mile further}}.
\]

\item A tin holds \(4\frac12\) litres of paint. Five jugs of \(\frac34\) of a litre are poured out. How much paint is left?
\textbf{Worked solution.}
Poured out:
\[
5\times\frac34=\frac{15}{4}=3\frac34.
\]
Paint left:
\[
4\frac12-3\frac34=\frac92-\frac{15}{4}
=\frac{18}{4}-\frac{15}{4}
=\frac34.
\]
So
\[
\boxed{\frac34\text{ litre left}}.
\]

\item A wall is \(7\frac12\) m wide. Panels are \(1\frac14\) m wide. How many panels fit across the wall?
\textbf{Worked solution.}
\[
7\frac12\div1\frac14=\frac{15}{2}\div\frac54
=\frac{15}{2}\times\frac45
=\frac{60}{10}
=6.
\]
So
\[
\boxed{6\text{ panels}}.
\]

\item A bottle holds \(1\frac13\) litres of water. Jo drinks \(\frac38\) of it. How much water does Jo drink?
\textbf{Worked solution.}
\[
1\frac13\times\frac38=\frac43\times\frac38
=\frac{12}{24}
=\frac12.
\]
So Jo drinks
\[
\boxed{\frac12\text{ litre}}.
\]

\item A journey takes \(2\frac14\) hours. Maria has driven \(\frac23\) of the way. How long is left, in minutes?
\textbf{Worked solution.}
Total time:
\[
2\frac14\text{ h}=\frac94\text{ h}=135\text{ min}.
\]
Maria has driven \(\frac23\), so the fraction remaining is
\[
1-\frac23=\frac13.
\]
Time left:
\[
\frac13\times135=45\text{ min}.
\]
So
\[
\boxed{45\text{ minutes}}.
\]

\item A rope \(12\frac12\) m long is cut into pieces \(1\frac34\) m long. How many whole pieces are there, and how much rope is left over?
\textbf{Worked solution.}
\[
12\frac12\div1\frac34=\frac{25}{2}\div\frac74
=\frac{25}{2}\times\frac47
=\frac{100}{14}
=\frac{50}{7}
=7\frac17.
\]
So there are \(7\) whole pieces.

Length of \(7\) pieces:
\[
7\times1\frac34=7\times\frac74=\frac{49}{4}=12\frac14.
\]
Left over:
\[
12\frac12-12\frac14=\frac14.
\]
Hence
\[
\boxed{7\text{ whole pieces, with }\frac14\text{ m left over}}.
\]

\item A recipe for \(4\) people uses \(2\frac12\) cups of rice and \(1\frac34\) litres of stock. How much rice and stock are needed for \(10\) people?
\textbf{Worked solution.}
Scale factor:
\[
\frac{10}{4}=\frac52.
\]
Rice:
\[
2\frac12\times\frac52=\frac52\times\frac52=\frac{25}{4}=6\frac14.
\]
Stock:
\[
1\frac34\times\frac52=\frac74\times\frac52=\frac{35}{8}=4\frac38.
\]
So for \(10\) people,
\[
\boxed{6\frac14\text{ cups of rice}}\quad\text{and}\quad
\boxed{4\frac38\text{ litres of stock}}.
\]
\end{enumerate}

%================================================================
\section*{Recurring decimals to fractions}

\begin{enumerate}[label=\arabic*., start=1, leftmargin=*, itemsep=6pt]
\item Use an algebraic method to write \(0.\dot3\) as a fraction.
\textbf{Worked solution.}
Let
\[
x=0.\dot3=0.333\ldots
\]
Then
\[
10x=3.333\ldots
\]
Subtract:
\[
10x-x=3.333\ldots-0.333\ldots
\]
so
\[
9x=3 \implies x=\frac13.
\]

\item Use an algebraic method to write \(0.\dot7\) as a fraction.
\textbf{Worked solution.}
Let
\[
x=0.\dot7=0.777\ldots
\]
Then
\[
10x=7.777\ldots
\]
Subtract:
\[
9x=7 \implies x=\frac79.
\]

\item Use an algebraic method to write \(0.\dot4\dot5\) as a fraction.
\textbf{Worked solution.}
Let
\[
x=0.\dot4\dot5=0.454545\ldots
\]
Then
\[
100x=45.454545\ldots
\]
Subtract:
\[
99x=45 \implies x=\frac{45}{99}=\frac{5}{11}.
\]

\item Use an algebraic method to write \(0.\dot1\dot2\) as a fraction in simplest form.
\textbf{Worked solution.}
Let
\[
x=0.\dot1\dot2=0.121212\ldots
\]
Then
\[
100x=12.121212\ldots
\]
Subtract:
\[
99x=12 \implies x=\frac{12}{99}=\frac{4}{33}.
\]

\item Use an algebraic method to write \(0.\dot2\dot7\) as a fraction in simplest form.
\textbf{Worked solution.}
Let
\[
x=0.\dot2\dot7=0.272727\ldots
\]
Then
\[
100x=27.272727\ldots
\]
Subtract:
\[
99x=27 \implies x=\frac{27}{99}=\frac{3}{11}.
\]

\item Use an algebraic method to write \(0.\dot1 2\dot3\) as a fraction in simplest form.
\textbf{Worked solution.}
Let
\[
x=0.\dot1 2\dot3=0.123123123\ldots
\]
Then
\[
1000x=123.123123123\ldots
\]
Subtract:
\[
999x=123 \implies x=\frac{123}{999}=\frac{41}{333}.
\]

\item Use an algebraic method to write \(0.1\dot6\) as a fraction in simplest form.
\textbf{Worked solution.}
Let
\[
x=0.1\dot6=0.1666\ldots
\]
Then
\[
10x=1.666\ldots
\]
and
\[
100x=16.666\ldots
\]
Subtract:
\[
90x=15 \implies x=\frac{15}{90}=\frac16.
\]

\item Use an algebraic method to write \(0.8\dot3\) as a fraction in simplest form.
\textbf{Worked solution.}
Let
\[
x=0.8\dot3=0.8333\ldots
\]
Then
\[
10x=8.333\ldots
\]
and
\[
100x=83.333\ldots
\]
Subtract:
\[
90x=75 \implies x=\frac{75}{90}=\frac56.
\]

\item Use an algebraic method to write \(1.\dot2\dot7\) as a mixed number in simplest form.
\textbf{Worked solution.}
Let
\[
x=1.\dot2\dot7=1.272727\ldots
\]
Then
\[
100x=127.272727\ldots
\]
Subtract:
\[
99x=126 \implies x=\frac{126}{99}=\frac{14}{11}=1\frac{3}{11}.
\]

\item Use an algebraic method to show that \(0.\dot9=1\).
\textbf{Worked solution.}
Let
\[
x=0.\dot9=0.999\ldots
\]
Then
\[
10x=9.999\ldots
\]
Subtract:
\[
10x-x=9.999\ldots-0.999\ldots
\]
so
\[
9x=9 \implies x=1.
\]
Hence
\[
0.\dot9=1.
\]
\end{enumerate}

\end{document}